\input{sol_pre.tex}
\input{q_2026-03-01_2.tex}
\begin{document}
\input{sol_head.tex}

%================= 1
\Savol{1}
12 ga böliniş $\iff$ 3 ga va 4 ga böliniş.\\
3 ga: $8+9+6+2+a+b=25+a+b$ $\Rightarrow$ $a+b\equiv2\pmod3$. \ 4 ga: $\overline{ab}$ 4 ga bölinadi.\\
Mos $\overline{ab}$: $08,20,32,44,56,68,80,92$. \ Köpaytmalar: $0,0,6,16,30,48,0,18$. Eng kattasi $6\cdot8=48$.
\Javob{B) 48}

%================= 2
\Savol{2}
$[\lg5]=0$ ($0<\lg5<1$), \ $\{\lg0{,}1\}=\{-1\}=0$, \ $[\ln3]=1$ ($e<3<e^2$).
\[\{3\cdot0+5\cdot0-7\cdot1\}=\{-7\}=0.\]
\Javob{A) 0}

%================= 3
\Savol{3}
Piyoda tezligi $v$, velosipedçiniki $4v$, masofa $S$. Yarim yölni piyoda $\frac S{2v}$, velosipedçi $\frac S{8v}$ soatda ötadi:
\[\frac S{2v}-\frac S{8v}=3\ \Rightarrow\ \frac{3S}{8v}=3\ \Rightarrow\ \frac Sv=8.\]
\Javob{B) 8}

%================= 4
\Savol{4}
Köylak $x$, şim $y$: \ $0{,}4x+0{,}45y=215\,000$, \ $0{,}45x+0{,}4y=210\,000$.\\
Qöşsak: $0{,}85(x+y)=425\,000\Rightarrow x+y=500\,000$; \ ayirsak: $0{,}05(y-x)=5000\Rightarrow y-x=100\,000$.\\
Demak $x=200\,000$.
\Javob{B) 200\,000}

%================= 5
\Savol{5}
Içkaridan boşlaymiz: $(4-1)^{-1}=\frac13$; \ $3+\frac13=\frac{10}3$, \ $\left(\frac{10}3\right)^{-1}=\frac3{10}$;\\
$2-\frac3{10}=\frac{17}{10}$, \ $\left(\frac{17}{10}\right)^{-1}=\frac{10}{17}$; \ $1+\frac{10}{17}=\frac{27}{17}$, \ $\left(\frac{27}{17}\right)^{-1}=\frac{17}{27}$.
\Javob{C) $\frac{17}{27}$}

%================= 6
\Savol{6}
$\sqrt[3]{6^2\sqrt6}=6^{\frac{2+\frac12}3}=6^{\frac56}$, \quad $\sqrt{6\sqrt[3]{6^2}}=6^{\frac{1+\frac23}2}=6^{\frac56}$.
\[\left((9{,}4-3{,}4)\cdot6^{\frac56}\right)^{\frac6{11}}=\left(6^{\frac{11}6}\right)^{\frac6{11}}=6.\]
\Javob{B) 6}

%================= 7
\Savol{7}
$(2\pm\sqrt3)^3=8\pm12\sqrt3+18\pm3\sqrt3=26\pm15\sqrt3$:
\[(2+\sqrt3)+(2-\sqrt3)=4.\]
\Javob{C) 4}

%================= 8
\Savol{8}
$a_2=2a_1\Rightarrow d=a_1$; \ $a_{10}=a_1+9d=10a_1=20\Rightarrow a_1=d=2$.\\
$S_6-S_5=a_6=a_1+5d=12$.
\Javob{B) 12}

%================= 9
\Savol{9}
$b_1=a_1+d$, $b_2=a_1+8d$, $b_3=a_1+43d$; \ $b_2^2=b_1b_3$:
\[a_1^2+16a_1d+64d^2=a_1^2+44a_1d+43d^2\ \Rightarrow\ 21d^2=28a_1d\ \Rightarrow\ a_1=\tfrac34d.\]
$b_1=\frac74d$, $b_2=\frac{35}4d$, $b_3=\frac{175}4d$ $\Rightarrow$ $\frac{217}4d=217$, \ $d=4$, $a_1=3$.\\
$S_n=\dfrac{6+4(n-1)}2\,n=n(2n+1)=820\Rightarrow n=20$ \ ($20\cdot41=820$).
\Javob{B) 20}

%================= 10
\Savol{10}
$x=y+2$: \quad $\dfrac{2(y+2)+4y+2}{2y+2}=\dfrac{6(y+1)}{2(y+1)}=3$.
\Javob{C) 3}

%================= 11
\Savol{11}
$\dfrac1{a+\frac1{b+\frac1c}}=\dfrac{bc+1}{abc+a+c}$, \quad $\dfrac1{a+\frac1b}=\dfrac b{ab+1}$.
\[\frac{bc+1}{abc+a+c}\cdot\frac{ab+1}b-\frac1{b(abc+a+c)}=\frac{ab^2c+ab+bc+1-1}{b(abc+a+c)}=\frac{b(abc+a+c)}{b(abc+a+c)}=1.\]
\Javob{D) 1}

%================= 12
\Savol{12}
$\sin11^\circ-\sin49^\circ=2\cos30^\circ\sin(-19^\circ)=-\sqrt3\sin19^\circ$;\\
$1-2\cos^254{,}5^\circ=-\cos109^\circ=\cos71^\circ=\sin19^\circ$.
\[\frac{-\sqrt3\sin19^\circ}{\sin19^\circ}=-\sqrt3.\]
\Javob{B) $-\sqrt3$}

%================= 13
\Savol{13}
$\sin^2x=\frac{1-\cos2x}2$, \ $\sin^22x=1-\cos^22x$, \ $t=\cos2x$:
\[\tfrac{\sqrt3}2-\tfrac{\sqrt3}2t+1-t^2=\tfrac{\sqrt3}2-\tfrac12-\sqrt3t\ \Rightarrow\ t^2-\tfrac{\sqrt3}2t-\tfrac32=0\ \Rightarrow\ t=\sqrt3\ (\text{yaroqsiz}),\ t=-\tfrac{\sqrt3}2.\]
$2x=\pm\frac{5\pi}6+2\pi k\Rightarrow x=\pm\frac{5\pi}{12}+\pi k$. \ $\left[-\frac\pi2;\pi\right]$ da: $-\frac{5\pi}{12},\ \frac{5\pi}{12},\ \frac{7\pi}{12}$ — 3 ta.
\Javob{B) 3}

%================= 14
\Savol{14}
$\dfrac{x^2+x-8}{x+1}=x-\dfrac8{x+1}=e$ deb olsak, \ $\left(\frac34\right)^e>\frac35\left(\frac54\right)^e\iff\left(\frac35\right)^e>\left(\frac35\right)^1\iff e<1$:
\[\frac{x^2+x-8-x-1}{x+1}<0\ \iff\ \frac{(x-3)(x+3)}{x+1}<0\ \iff\ x\in(-\infty;-3)\cup(-1;3).\]
\Javob{B) $(-\infty;-3)\cup(-1;3)$}

%================= 15
\Savol{15}
$\log_{0{,}5}u>0\iff0<u<1$: \ $0<\log_2\log_{x-1}9<1\iff1<\log_{x-1}9<2$.\\
Bu faqat $x-1>1$ da böladi: $x-1<9$ va $(x-1)^2>9$ $\Rightarrow$ $4<x<10$.\\
Butunlar: $5,6,7,8,9$ — 5 ta.
\Javob{B) 5}

%================= 16
\Savol{16}
$(x^2-5x-6)+(7+6x-x^2)=x+1$. \ $|a|+|b|=a+b\iff a\ge0,\ b\ge0$:
\[x^2-5x-6\ge0\Rightarrow x\le-1\ \text{yoki}\ x\ge6;\qquad -x^2+6x+7\ge0\Rightarrow-1\le x\le7.\]
Yeçim: $x=-1$ va $[6;7]$ — çeksiz köp.
\Javob{D) çeksiz köp yeçimga ega}

%================= 17
\Savol{17}
$f(x)=8x^3-6x^2+3x-1$, \ $f'(x)=24x^2-12x+3$, \ $D=144-288<0$ $\Rightarrow$ $f'>0$, $f$ ösuvçi.\\
Kubik funksiya ösuvçi bölsa, bitta ildizga ega.
\Javob{A) 1}

%================= 18
\Savol{18}
Umumiy maxraj: \ surat $=(x+2)(x+3)+2(x+1)(x+2)-3(x+1)(x+3)=1-x$.
\[\frac{1-x}{(x+1)(x+2)(x+3)}>0\iff\frac{x-1}{(x+1)(x+2)(x+3)}<0\ \Rightarrow\ x\in(-3;-2)\cup(-1;1).\]
\Javob{C) $(-3;-2)\cup(-1;1)$}

%================= 19
\Savol{19}
$A^2<B^2\iff(A-B)(A+B)<0$:
\[(x^2-6x-6)(3x^2+4x+2)<0.\]
$3x^2+4x+2>0$ ($D<0$) $\Rightarrow$ $x^2-6x-6<0$, \ $3-\sqrt{15}<x<3+\sqrt{15}$ \ ($\approx-0{,}87<x<6{,}87$).\\
Butunlar: $0,1,\dots,6$, yiğindisi $21$.
\Javob{B) 21}

%================= 20
\Savol{20}
Nollari $-2;0;2$ $\Rightarrow$ $f(x)=ax(x^2-4)$. \ $f(-1)=3a=3\Rightarrow a=1$.
\[f(-3)=-3\cdot(9-4)=-15.\]
\Javob{C) $-15$}

%================= 21
\Savol{21}
$u=3x+1$, \ $x=\frac{u-1}3$, \ $dx=\frac{du}3$:
\[\int\frac{u-1}{9\sqrt u}\,du=\frac19\left(\frac23u^{\frac32}-2u^{\frac12}\right)+C=\frac29\left(\frac13(3x+1)^{\frac32}-(3x+1)^{\frac12}\right)+C.\]
\Javob{D}

%================= 22
\Savol{22}
$x=\frac\pi4$ atrofida $\log_{\sin x}\cos x\approx1>\ln x$ $\Rightarrow$ $f(x)=\dfrac{\ln\cos x}{\ln\sin x}-\ln x$.\\
$x=\frac\pi4$ da $L=\ln\sin x=\ln\cos x=-\frac{\ln2}2$:
\[f'\!\left(\tfrac\pi4\right)=\frac{-\tg x\cdot L-L\ctg x}{L^2}-\frac4\pi=\frac{-2L}{L^2}-\frac4\pi=-\frac2L-\frac4\pi=\frac4{\ln2}-\frac4\pi.\]
\Javob{B) $-\frac4\pi+\frac4{\ln2}$}

%================= 23
\Savol{23}
$3x-4$ va $2x+1$ butun $\Rightarrow$ $x=(3x-4)-(2x+1)+5$ butun.\\
Şartlar: $2x+1\le8$, \ $3x-4\le8$, \ $(3x-4)+(2x+1)>8$ $\Rightarrow$ $2{,}2<x\le3{,}5$ $\Rightarrow$ $x=3$.\\
Tomonlar: $5,7,8$; \ $P=20$.
\Javob{D) 20}

%================= 24
\Savol{24}
I. $y_1(-x)=f(x)\cdot(-g(x))=-y_1(x)$ — toq. \ III. $y_3(-x)=\frac{f(x)}{-g(x)}=-y_3(x)$ — toq.\\
II. $y_2(-x)=f(x)-g(x)$ — umuman olganda na juft, na toq.
\Javob{C) 2}

%================= 25
\Savol{25}
Aylanalar radiusi $r$. Ular özaro urinadi: $O_1O_2O_3$ — tomoni $2r$ muntazam uçburçak.\\
$O_3$ markaz $A$ burçak bissektrisasida, $AC$ dan $r$ masofada: $AK=r\ctg30^\circ=\sqrt3r$.
\fig{\begin{tikzpicture}[scale=0.62]
\coordinate (A) at (0,0); \coordinate (C) at (5.464,0); \coordinate (B) at (2.732,4.732);
\draw (A)--(B)--(C)--cycle;
\foreach \p in {(1.732,1),(3.732,1),(2.732,2.732)} \draw \p circle (1);
\draw (1.732,1)--(3.732,1)--(2.732,2.732)--cycle;
\draw[ax] (A)--(1.732,1)--(1.732,0); \draw[ax] (3.732,1)--(3.732,0);
\node[below left] at (A){$A$}; \node[below right] at (C){$C$}; \node[above] at (B){$B$};
\node[below,font=\scriptsize,aux] at (0.87,0){$\sqrt3r$}; \node[below,font=\scriptsize,aux] at (2.73,0){$2r$};
\node[below,font=\scriptsize,aux] at (4.6,0){$\sqrt3r$}; \node[font=\scriptsize] at (1.4,1.25){$O_3$};
\end{tikzpicture}}
$AC=2r+2\sqrt3r$. \ Öxşaş uçburçaklar: \ $\dfrac{S_{ABC}}{S_{O_1O_2O_3}}=\left(\dfrac{2r+2\sqrt3r}{2r}\right)^2=(1+\sqrt3)^2=4+2\sqrt3$.
\Javob{B) $4+2\sqrt3$}

%================= 26
\Savol{26}
$BH$ — balandlik, $AH=32$, katta aylana markazi $I$ ($IH=16$). $AI$ — bissektrisa: $\tg\frac A2=\frac{16}{32}=\frac12$,
\[\tg A=\frac{2\cdot\frac12}{1-\frac14}=\frac43\ \Rightarrow\ BH=32\cdot\frac43=\frac{128}3.\]
Katta aylanaga $AC$ ga parallel urinma (qizil) $B$ dan $\frac{128}3-32=\frac{32}3$ masofada. Kiçik aylana — kesilgan uçburçakning içki aylanasi, u $ABC$ ga öxşaş, koeffitsiyent $\frac{32/3}{128/3}=\frac14$:
\fig{\begin{tikzpicture}[scale=0.055]
\coordinate (A) at (-32,0); \coordinate (C) at (32,0); \coordinate (B) at (0,42.67);
\draw (A)--(B)--(C)--cycle; \draw (0,16) circle (16); \draw (0,36) circle (4);
\draw[ax] (-8,32)--(8,32); \draw[ax] (B)--(0,0); \draw[ax] (A)--(0,16);
\node[below left] at (A){$A$}; \node[below right] at (C){$C$}; \node[above] at (B){$B$};
\node[below] at (0,0){$H$}; \node[right,font=\scriptsize] at (0,16){$I$};
\end{tikzpicture}}
\[r=16\cdot\frac14=4.\]
\Javob{B) 4}

%================= 27
\Savol{27}
$S_{ABD}=S_{ACD}$ (asos $AD$, balandlik bir xil) $\Rightarrow$ $S_{AOB}=S_{COD}=S$.\\
$\dfrac{S_{AOB}}{S_{BOC}}=\dfrac{AO}{OC}=\dfrac{S_{AOD}}{S_{COD}}$ $\Rightarrow$ $S^2=S_{BOC}\cdot S_{AOD}=2\cdot8=16$, \ $S=4$.
\Javob{A) 4}

%================= 28
\Savol{28}
$CP=CB=CD=CE$ (kvadrat va muntazam uçburçaklar tomonlari).\\
$\angle PCE=\angle PCD+\angle DCE=(90^\circ-60^\circ)+60^\circ=90^\circ$ $\Rightarrow$ $\triangle PCE$ teng yonli töğri burçakli, $\angle CEP=45^\circ$.\\
$\triangle KDE$ da: $\angle KDE=60^\circ$, \ $\angle KED=\angle CED-\angle CEP=60^\circ-45^\circ=15^\circ$:
\[\angle DKE=180^\circ-60^\circ-15^\circ=105^\circ.\]
\Javob{D) $105^\circ$}

%================= 29
\Savol{29}
$D$ — $AC$ örtasi: $D(1;6;-1)$. \ $\overrightarrow{BD}=(3;-1;0)$, \ $\overrightarrow{BC}=(6;0;3)$.
\[\cos\angle CBD=\frac{18}{\sqrt{10}\cdot\sqrt{45}}=\frac{3\sqrt2}5,\qquad \sin\angle CBD=\sqrt{1-\frac{18}{25}}=\frac{\sqrt7}5.\]
\Javob{A) $\frac{\sqrt7}5$}

%================= 30
\Savol{30}
$\angle C=90^\circ$, $AC=BC=a$, $AH\perp\alpha$. \ $\angle ACH=45^\circ\Rightarrow AH=\frac a{\sqrt2}$; \ $AB=a\sqrt2$.
\fig{\begin{tikzpicture}[scale=0.8]
\draw (0,0)--(5,0)--(6,1.6)--(1,1.6)--cycle; \node at (5.4,0.25){$\alpha$};
\coordinate (C) at (1.6,0.5); \coordinate (B) at (4.6,0.6); \coordinate (H) at (2.6,1.2); \coordinate (A) at (2.6,2.6);
\draw (A)--(C)--(B)--cycle; \draw[ax] (A)--(H); \draw[ax] (C)--(H)--(B);
\node[left] at (C){$C$}; \node[right] at (B){$B$}; \node[above] at (A){$A$}; \node[below right] at (H){$H$};
\end{tikzpicture}}
\[\sin\angle ABH=\frac{AH}{AB}=\frac{a/\sqrt2}{a\sqrt2}=\frac12\ \Rightarrow\ \angle ABH=30^\circ.\]
\Javob{A) $30^\circ$}

%================= 31
\Savol{31}
Kamida bitta tögarakka: $25-10=15$. \ Ikkalasiga: $9+8-15=2$.
\Javob{A) 2}

%================= 32
\Savol{32}
Uçburçakning ikki uçi bir töğri çiziqda, uçinçisi boşqasida böladi:
\[C_3^2\cdot4+C_4^2\cdot3=12+18=30.\]
\Javob{D) 30}

%================= 33-35
\newpage
\Savol{33}
Prizmaga şar içki çizilgan $\Rightarrow$ asosga aylana içki çiziladi va prizma balandligi $H=2r$.\\
Trapetsiya: $AB+CD=8+18\Rightarrow$ yon tomon $13$; \ balandlik $h=\sqrt{13^2-5^2}=12$, \ $H=h=12$.
\fig{\begin{tikzpicture}[scale=0.28]
\coordinate (A) at (0,0); \coordinate (D) at (18,0); \coordinate (C) at (13,12); \coordinate (B) at (5,12);
\draw (A)--(B)--(C)--(D)--cycle; \draw (9,6) circle (6);
\draw[axs] (A)--(C); \draw[ax] (B)--(D); \draw[ax] (2.5,6)--(15.5,6);
\draw[axs] (4.6,5.3)--(13.4,5.3);
\node[below left] at (A){$A$}; \node[below right] at (D){$D$}; \node[above right] at (C){$C$}; \node[above left] at (B){$B$};
\node[below] at (9,0){$18$}; \node[above] at (9,12){$8$};
\end{tikzpicture}}
\textbf{33.} $AC^2=(5+8)^2+12^2=313$ $\Rightarrow$ $S=AC\cdot H=12\sqrt{313}$.\\
\textbf{34.} Diagonallar kesişgan nuqtadan asoslarga parallel kesma: $\dfrac{2\cdot8\cdot18}{8+18}=\dfrac{144}{13}$; \ $S=\dfrac{144}{13}\cdot12=\dfrac{1728}{13}=132\dfrac{12}{13}$.\\
\textbf{35.} Örta çiziq $\frac{8+18}2=13$; \ $S=13\cdot12=156$.
\Javob{33 — A; \quad 34 — B; \quad 35 — E}

%================= 36
\newpage
\Savol{36}
Viyet: $r+s+t=-a$, \ $rs+st+tr=b$, \ $rst=-1$.\\
$g$ uçun: \ $r^2+s^2+t^2=-m$, \ $r^2s^2+s^2t^2+t^2r^2=n$, \ $r^2s^2t^2=-p$.
\[-m=(r+s+t)^2-2(rs+st+tr)=a^2-2b,\qquad p=-(rst)^2=-1,\]
\[n=(rs+st+tr)^2-2rst(r+s+t)=b^2-2(-1)(-a)=b^2-2a.\]
$g(-1)=-1+m-n+p=-5$ $\Rightarrow$ $m-n=-3$, ya'ni $n-m=3$.

a) $n-m-p=3-(-1)=4$.

b) $n-m=b^2-2a+a^2-2b=3$ $\Rightarrow$ $(a^2-2a+1)+(b^2-2b+1)=5$, \ $(a-1)^2+(b-1)^2=5$.
\Javob{a) 4; \quad b) 5}

%================= 37
\newpage
\Savol{37}
$\sin^22x=4\sin^2x\cos^2x$. \ $\cos x=0$ bölsa, $\sin^4x=1\ne0$ — ildiz emas. $\cos^4x$ ga bölamiz, $t=\tg^2x$:
\[t^2+8t-9=0\ \Rightarrow\ t=1\ \ (t=-9\ \text{yaroqsiz})\ \Rightarrow\ \tg x=\pm1,\quad x=\frac\pi4+\frac{\pi k}2.\]
a) $-\frac\pi2\le\frac\pi4+\frac{\pi k}2\le\frac{5\pi}2$ $\Rightarrow$ $-1{,}5\le k\le4{,}5$, \ $k=-1,0,1,2,3,4$ — 6 ta.

b) $[-\pi;\pi]$ da: $-\frac{3\pi}4,\ -\frac\pi4,\ \frac\pi4,\ \frac{3\pi}4$; \ yiğindisi $0$.
\Javob{a) 6; \quad b) 0}

%================= 38
\newpage
\Savol{38}
$f(x)=ax^2+bx+2$.

a) $f(x+5)-f(x)=a(10x+25)+5b=x+1$ $\Rightarrow$ $10a=1$, $25a+5b=1$ $\Rightarrow$ $a=\frac1{10}$, $b=-\frac3{10}$.
\[f(2025)=\frac{2025^2-3\cdot2025}{10}+2=\frac{2025\cdot2022}{10}+2=409\,455+2=409\,457.\]

b) $f(x-9)-f(x)=a(-18x+81)-9b=x-3$ $\Rightarrow$ $-18a=1$, $81a-9b=-3$ $\Rightarrow$ $a=-\frac1{18}$, $b=-\frac16$.
\[f(2025)=-\frac{2025^2}{18}-\frac{2025}6+2=-\frac{2025(2025+3)}{18}+2=-\frac{2025\cdot2028}{18}+2=-228\,150+2=-228\,148.\]
\Javob{a) 409\,457; \quad b) $-228\,148$}

%================= 39
\newpage
\Savol{39}
a) Urinma $y=2x+1$ bilan kesişmaydi $\Rightarrow$ unga parallel: $f'(x_0)=2x_0-6=2$ $\Rightarrow$ $x_0=4$, \ $y_0=16-24+12=4$.

b) $x^2-6x+12=-2x+2\iff x^2-4x+10=0$, $D<0$ — parabola töğri çiziqni kesmaydi. Eng yaqin nuqta — urinmasi $y=-2x+2$ ga parallel bölgan nuqta:
\[2x_0-6=-2\Rightarrow x_0=2,\ y_0=4;\qquad d=\frac{|2\cdot2+4-2|}{\sqrt{2^2+1^2}}=\frac6{\sqrt5}=\frac{6\sqrt5}5.\]
\fig{\begin{tikzpicture}[scale=0.45]
\draw[->] (-1.5,0)--(6.5,0) node[below]{$x$}; \draw[->] (0,-1.5)--(0,8.5) node[left]{$y$};
\draw[thick,domain=0.6:5.4,samples=50] plot (\x,{\x*\x-6*\x+12});
\draw[domain=-1:1.6] plot (\x,{-2*\x+2});
\draw[ax,domain=0.2:4] plot (\x,{-2*\x+8});
\fill (2,4) circle (2.5pt) node[right]{$A(2;4)$};
\draw[axs] (2,4)--(-0.4,2.8);
\node[font=\scriptsize,right] at (1.6,-1.2){$y=-2x+2$};
\end{tikzpicture}}
\Javob{a) 4; \quad b) $\frac{6\sqrt5}5$}

%================= 40
\newpage
\Savol{40}
Kesişiş: $x^2=\sqrt x\Rightarrow x=0,\ x=1$. \ $[0;1]$ da $\sqrt x\ge x^2$.
\fig{\begin{tikzpicture}[scale=2]
\fill[okgreen!30,domain=0:1,samples=40] plot (\x,{sqrt(\x)}) -- plot[domain=1:0] (\x,{\x*\x}) -- cycle;
\draw[->] (-0.1,0)--(1.4,0) node[below]{$x$}; \draw[->] (0,-0.1)--(0,1.3) node[left]{$y$};
\draw[thick,domain=0:1.15,samples=40] plot (\x,{\x*\x}) node[above]{$y=x^2$};
\draw[thick,domain=0:1.3,samples=60] plot (\x,{sqrt(\x)}) node[right]{$y=\sqrt x$};
\node[below] at (1,0){$1$}; \draw[ax] (1,0)--(1,1);
\end{tikzpicture}}
a) $S=\displaystyle\int_0^1\left(\sqrt x-x^2\right)dx=\frac23-\frac13=\frac13$.

b) $V=\pi\displaystyle\int_0^1\left((\sqrt x)^2-(x^2)^2\right)dx=\pi\left(\frac12-\frac15\right)=\frac{3\pi}{10}\approx\frac{3\cdot3}{10}=0{,}9$.
\Javob{a) $\frac13$; \quad b) 0,9}

%================= 41
\newpage
\Savol{41}
$\angle A=\alpha$, $\angle B=\beta$, $\angle C=\gamma$. $E$ va $F$ — taşqi aylanada yotadi ($B$ bissektrisasi va $AC$ ning örta perpendikulyari $AC$ yoyining örtasida kesişadi; $F$ ham şunday).
\fig{\begin{tikzpicture}[scale=1.25]
\coordinate (O) at (0,0); \coordinate (A) at (-1.5,-0.866); \coordinate (C) at (1.5,-0.866); \coordinate (B) at (-0.592,1.628);
\coordinate (F) at (1.327,1.113); \coordinate (E) at (0,-1.732);
\draw[ax] (O) circle (1.732);
\draw (A)--(B)--(C)--cycle; \draw (A)--(F); \draw (B)--(E); \draw (E)--(O)--(F);
\draw[axs] (O)--(A); \draw[axs] (O)--(B);
\fill (O) circle (1pt);
\node[left] at (A){$A$}; \node[right] at (C){$C$}; \node[above] at (B){$B$}; \node[right] at (F){$F$}; \node[below] at (E){$E$}; \node[above left] at (O){$O$};
\end{tikzpicture}}
$OA=OF=R$ $\Rightarrow$ $\angle OAF=\angle AFO=5^\circ$. \ $\angle AOC=2\beta$ $\Rightarrow$ $\angle OAC=90^\circ-\beta$, \ $\angle FAC=\frac\alpha2$:
\[\angle OAF=\frac\alpha2-(90^\circ-\beta)=\frac{\alpha+2\beta-180^\circ}2=\frac{\beta-\gamma}2=5^\circ\ \Rightarrow\ \beta-\gamma=10^\circ.\]
Xuddi şunday $\angle OBE=\angle BEO=10^\circ$: \ $\angle OBE=\frac\beta2-(90^\circ-\alpha)=\frac{\alpha-\gamma}2=10^\circ\Rightarrow\alpha-\gamma=20^\circ$.

a) $(\gamma+20^\circ)+(\gamma+10^\circ)+\gamma=180^\circ\Rightarrow\gamma=50^\circ$ \ ($\beta=60^\circ$, $\alpha=70^\circ$).

b) $OF=R=\dfrac{AC}{2\sin\beta}=\dfrac3{2\cdot\frac{\sqrt3}2}=\sqrt3$.
\Javob{a) $50^\circ$; \quad b) $\sqrt3$}

%================= 42
\newpage
\Savol{42}
Trapetsiya aylanaga içki çizilgan $\Rightarrow$ teng yonli, $\smile AB=\smile CD$. Urinma $DC$ ga parallel $\Rightarrow$ $OE\perp DC$, ya'ni $E$ — $CD$ yoyining örtasi; $EF\parallel BC$ bölgani uçun $F$ — $AB$ yoyining örtasi.
\fig{\begin{tikzpicture}[scale=0.055]
\draw (0,0) circle (37.04);
\coordinate (A) at (-37,-1.73); \coordinate (D) at (37,-1.73); \coordinate (B) at (-28,24.25); \coordinate (C) at (28,24.25);
\coordinate (E) at (35,12.12); \coordinate (F) at (-35,12.12); \coordinate (O) at (0,0);
\draw (A)--(B)--(C)--(D)--cycle; \draw (F)--(E);
\draw[axs] (A)--(C); \draw[ax] (O)--(E); \draw[ax] (C)--(28,-1.73);
\fill (O) circle (1.2);
\node[left] at (A){$A$}; \node[right] at (D){$D$}; \node[above left] at (B){$B$}; \node[above right] at (C){$C$};
\node[right] at (E){$E$}; \node[left] at (F){$F$}; \node[below] at (O){$O$}; \node[below] at (28,-1.73){$H$};
\node[above] at (0,24.25){$56$};
\end{tikzpicture}}
a) $\smile FE=\smile FB+\smile BC+\smile CE=\frac12\smile AB+\smile BC+\frac12\smile CD=\smile AB+\smile BC=\smile AC$.\\
Teng yoylarni teng vatarlar tortadi: \ $AC=EF=70$.

b) $O$ dan $BC$ gaça: $\sqrt{R^2-28^2}=\sqrt{1372-784}=14\sqrt3$. \ $AD=2x$, $O$ dan $AD$ gaça $d=\sqrt{1372-x^2}$, $h=14\sqrt3+d$.\\
$CH\perp AD$ da $AH=\frac{AD+BC}2=x+28$: \ $AC^2=h^2+(x+28)^2$:
\[(14\sqrt3+d)^2+(x+28)^2=4900\ \Rightarrow\ 588+28\sqrt3d+1372+56x+784=4900\ \Rightarrow\ \sqrt3d=77-2x.\]
$3(1372-x^2)=(77-2x)^2\Rightarrow x^2-44x+259=0\Rightarrow x=37$ \ ($x=7$ da $AD=14<56$ — yaroqsiz).\\
$d=\sqrt{1372-1369}=\sqrt3$, \ $h=15\sqrt3$, \ $AD=74$.
\[S=\frac{74+56}2\cdot15\sqrt3=975\sqrt3.\]
\Javob{a) 70; \quad b) $975\sqrt3$}

%================= 43
\newpage
\Savol{43}
Aylana radiusi $r$. U sakkizburçakka içki çizilgan $\Rightarrow$ markazdan sakkizburçak tomonigaça masofa $r$.\\
Oltiburçak aylanaga içki çizilgan, tomoni $r$ $\Rightarrow$ markazdan tomonigaça (apofema) $r\cos30^\circ=\frac{\sqrt3}2r$.
\fig{\begin{tikzpicture}[scale=1.5]
\draw (0,0) circle (1);
\draw (0:1)--(60:1)--(120:1)--(180:1)--(240:1)--(300:1)--cycle;
\draw (22.5:1.0824)--(67.5:1.0824)--(112.5:1.0824)--(157.5:1.0824)--(202.5:1.0824)--(247.5:1.0824)--(292.5:1.0824)--(337.5:1.0824)--cycle;
\draw[axs] (0,0)--(0,-1); \draw[ax] (0,0)--(30:0.866); \fill (0,0) circle (0.8pt);
\node[font=\scriptsize,aux,right] at (0,-0.55){$r$}; \node[font=\scriptsize,aux] at (0.55,0.5){$\frac{\sqrt3}2r$};
\end{tikzpicture}}
a) $\dfrac{\frac{\sqrt3}2r}r=\dfrac{\sqrt3}2$.

b) $S_6=6\cdot\frac{\sqrt3}4r^2=\frac{3\sqrt3}2r^2$. \ Sakkizburçak — 8 ta uçburçak, asosi $2r\tg22{,}5^\circ$, balandligi $r$:
\[S_8=8r^2\tg22{,}5^\circ=8(\sqrt2-1)r^2;\qquad \frac{S_6}{S_8}=\frac{3\sqrt3}{16(\sqrt2-1)}=\frac{3\sqrt3(\sqrt2+1)}{16}.\]
\Javob{a) $\frac{\sqrt3}2$; \quad b) $\frac{3\sqrt3(\sqrt2+1)}{16}$}

%================= 44
\newpage
\Savol{44}
Asoslar yuzlari nisbati 9 $\Rightarrow$ tomonlar $a=3b$ va $b$. Apofemalar orqali kesim — teng yonli trapetsiya (asoslari $3b$, $b$, yon tomoni $l$ — apofema), şar katta aylanasi unga içki çizilgan:
\fig{\begin{tikzpicture}[scale=0.38]
\coordinate (P) at (-5.196,0); \coordinate (Q) at (5.196,0); \coordinate (R) at (1.732,6); \coordinate (S) at (-1.732,6);
\draw (P)--(Q)--(R)--(S)--cycle; \draw (0,3) circle (3);
\draw[ax] (0,0)--(0,6); \draw[ax] (1.732,6)--(1.732,0);
\node[below] at (0,0){$3b$}; \node[above] at (0,6){$b$}; \node[right] at (3.6,3){$l$}; \node[right,font=\scriptsize,aux] at (0,1.5){$h$};
\end{tikzpicture}}
$2l=3b+b\Rightarrow l=2b$; \ $h=\sqrt{l^2-b^2}=\sqrt3\,b$.\\
Töla sirt: $9b^2+b^2+4\cdot\frac{3b+b}2\cdot2b=26b^2=312$ $\Rightarrow$ $b^2=12$, \ $h=\sqrt3\cdot2\sqrt3=6$.

a) $V=\dfrac h3\left(S_1+S_2+\sqrt{S_1S_2}\right)=\dfrac63(108+12+36)=312$.

b) $R=\frac h2=3$, \ $V_{sh}=\frac43\pi\cdot27=36\pi$.
\Javob{a) 312; \quad b) $36\pi$}

%================= 45
\newpage
\Savol{45}
Balandlik $h$, asos diagonalining yarmi $x$: \ $x^2+h^2=4$. Asos — diagonali $2x$ bölgan kvadrat: $S=2x^2$.
\fig{\begin{tikzpicture}[scale=0.9]
\coordinate (A) at (0,0); \coordinate (B) at (3,0); \coordinate (C) at (4,1.2); \coordinate (D) at (1,1.2);
\coordinate (O) at (2,0.6); \coordinate (S) at (2,3);
\draw (A)--(B)--(C); \draw[dashed] (C)--(D)--(A);
\draw (S)--(A); \draw (S)--(B); \draw (S)--(C); \draw[dashed] (S)--(D);
\draw[ax] (S)--(O); \draw[ax] (O)--(C); \draw[ax] (A)--(O);
\node[left] at (S){$S$}; \node[below] at (O){$O$};
\node[right,font=\scriptsize] at (3,2.1){$2$}; \node[right,font=\scriptsize,aux] at (2,1.8){$h$}; \node[below,font=\scriptsize,aux] at (3,0.9){$x$};
\end{tikzpicture}}
\[V(h)=\frac13\cdot2(4-h^2)\cdot h=\frac23(4h-h^3),\qquad V'(h)=\frac23(4-3h^2)=0\Rightarrow h=\frac2{\sqrt3}.\]
$h<\frac2{\sqrt3}$ da $V'>0$, keyin $V'<0$ — maksimum.

a) $x^2=4-\frac43=\frac83$, \ $S=2x^2=\frac{16}3$.

b) $V=\frac13\cdot\frac{16}3\cdot\frac2{\sqrt3}=\frac{32}{9\sqrt3}=\frac{32\sqrt3}{27}$.
\Javob{a) $\frac{16}3$; \quad b) $\frac{32\sqrt3}{27}$}

\ShowJavoblar
\end{document}
